\documentclass{article}
\usepackage{graphicx} 
\usepackage{xcolor}
\usepackage{amsmath}
\usepackage{amsfonts}
\usepackage{amssymb}
\usepackage{amsthm}
\usepackage[margin=1in]{geometry}
\usepackage{mathtools}
\title{Diseasenowcasting}
\author{Rodrigo Zepeda}
\date{August 2026}

\setcounter{tocdepth}{4}
\setcounter{secnumdepth}{4}

\newtheorem{theorem}{Theorem}
\newtheorem{proposition}{Proposition}
\newtheorem{corollary}{Corollary}
\newtheorem{lemma}{Lemma}
\theoremstyle{definition}
\newtheorem{definition}{Definition}
\theoremstyle{remark}
\newtheorem{remark}{Remark}

\newcommand{\Pois}{\operatorname{Poisson}}
\newcommand{\NB}{\operatorname{NB}}
\newcommand{\Sk}{\operatorname{Skellam}}
\newcommand{\E}{\operatorname{\mathbb{E}}}
\newcommand{\Var}{\operatorname{Var}}
\newcommand{\Cov}{\operatorname{Cov}}
\newcommand{\Prb}{\operatorname{\mathbb{P}}}

\begin{document}

\maketitle

\section{Methods}

Our methods aim to nowcast the number of events that have yet to be recorded due to reporting delays. As an example consider an individual with influenza-like symptoms. 
The event might correspond to symptom onset (the symptoms making it a \textit{suspected} case). The report date might correspond to the day the individual was seen by a medical professional for the first time and reported to the healthcare system. Finally the validation date might correspond to the date their influenza test results were registered in the same system. 

\begin{figure}[!htb]
    \centering    \includegraphics[width=0.7\linewidth]{diagram_confirmed.pdf}
    \caption{We consider the evolution of a suspected case from the event (e.g. symptom onset) to the report (e.g. registry as suspected into the healthcare system) to the validation (e.g. negative test).}
    \label{fig:placeholder}
\end{figure}

In general our models consider the workflow: event $\to$ report $\to$ validation. We  nowcast for scenarios where there is no validation as well as scenarios where the validation consists of only rejections (e.g. cases that lose eligibility), or only confirmations (e.g. systems that only keep the true positives). We also nowcast scenarios where only cumulative counts are available and the latest \textit{best estimate} of the total number of validated cases is reported.   

\subsection{Notation}
For each suspected case we consider the following notation:
\begin{itemize}
    \item \textbf{Now}, $\tau$, denotes the maximum observable time (\textit{i.e.} the most recent time where a case, a report, or a validation could have occurred). 
    \item \textbf{Event date}, $t$, denotes the date a suspected case occurred with $t \geq 0$. 
    \item \textbf{Report date}, $r$, denotes the date a suspected case was recorded with $r \geq t$. A case is reported by the now if and only if $r \leq \tau$. 
    \item \textbf{Validation date}, $v$ denotes the date a reported suspected case was either confirmed as a true case or retracted. It is bounded by the report with $v \geq r$. We allow for $v = \infty$ for a case that is never validated.   
    \item \textbf{Reporting delay}, $d^{\text{rpt}}$, is the difference between the reporting date and the event date (\textit{i.e.} $d^{\text{rpt}} = r - t$). The minimum delay possible is $d^{\text{rpt}} = 0$. 
    \item \textbf{Validation delay}, $d^{\text{val}}$, is the difference between the validation date and the report date (\textit{i.e.} $d^{\text{val}} = v - r$). The minimum delay possible is $d^{\text{val}} = 0$. For $v = \infty$ we write $d^{\text{val}} = \infty$.
    \item \textbf{Report's age,} $a$, denotes the difference between the now and the report date ($a = \tau - r$) for a report that is already observable ($r \leq \tau$).   
    \item \textbf{Maximum observable delay}, $d^{\star}_t$, denotes the maximum observable delay for an event that occurred at time $t$. Specifically, for a now $\tau$, we define it as $d^{\star}_t = \tau - t$.
    \item \textbf{Observed validation state}, $s$, is \texttt{confirmed} if the suspected case is observed to be a true case, \texttt{retracted} if it is observed not to be a true case, and \texttt{pending} if the validation has yet to be observed.
\end{itemize}
The nowcasting objective of the model is, at time $\tau \geq t$ (i.e. the \textit{now}), predict $N_t^+$ the overall number of true cases that happened at event time $t$ including those that have yet to be reported or whose validation is still pending. 


\subsubsection{Assumptions}
For a fixed event time $t$ we make the following assumptions:
\begin{enumerate}    
    \item \textbf{Temporal ordering:} 
    All reports happen after the event and all the validations happen after the report: $0 \leq t \leq r \leq v$. This follows the scheme: event $\to$ report $\to$ validation.   
    \item \textbf{Cases might never get validated:} The validation delay may take values in $\{0,1,2,\dots\} \cup \{\infty\}$.
    \item \textbf{Observed validation states:} A reported case observed at the now, $\tau$, can be \texttt{confirmed}, \texttt{retracted} or \texttt{pending}. A \texttt{pending} record is one whose eventual status is still unknown.
    \item \textbf{Parameters and Bayesian conditioning:} The vector $\theta_t$ contains all model parameters and has prior density $\pi(\theta_t)$. All probability distributions, probability masses, densities, expectations, and likelihood contributions below are understood to be conditional on $\theta_t$.
\end{enumerate}


\subsection{General process}

\subsubsection{Linelist data}
Let
\begin{equation}
    M_t \in \{0,1,2,\dots\}
\end{equation}
denote the \textbf{total number of events} that occurred at event time $t$. We refer to $\{M_t\}$ as the (latent) \textbf{epidemic process}. Conditional on $M_t = m$, let 
\begin{equation}
    X_{t,1}, X_{t,2}, \dots, X_{t,m}
\end{equation}
be i.i.d. marks with law $Q_t$ independent of $M_t$. A mark may contain as components, depending on the observation regime, the following:
\begin{equation}
X_{t,j} = (D^{\text{rpt}}_{t,j}, D^{\text{val}}_{t,j}, Y_{t,j}), 
\end{equation}
where $D^{\text{rpt}}$ is the \textbf{event-to-report delay}, $D^{\text{val}}$ the \textbf{report-to-validation delay}, and $Y$ is an \textbf{indicator with $Y = 1$ if it is a true case} and $0$ otherwise. We denote the probability of being a true positive as $p = \Prb(Y_{t,j} = 1\mid\theta_t)$. We refer to data where at least some of the records associated to each of the marks $X_{t,j}$ are available as line-list data. That is, datasets where each individual observation is available to the analyst. We discuss aggregated data in a later section.

Conditional on $\theta_t$, we assume the components of $X_{t,j}$ factorize as follows: the latent case indicator is Bernoulli, $Y_{t,j}\sim\operatorname{Bernoulli}(p)$; the reporting delay is independent of the latent case indicator, $D_{t,j}^{\mathrm{rpt}}\perp Y_{t,j}$; and the validation delay is conditionally independent of the reporting delay given the latent case indicator, $D_{t,j}^{\mathrm{val}}\perp D_{t,j}^{\mathrm{rpt}}\mid Y_{t,j}$. Thus, for a finite validation delay $\ell$ and $y\in\{0,1\}$,
\begin{equation}\label{markfactorization}
\begin{split}
\Prb\left(D_{t,j}^{\mathrm{rpt}}=d,\,Y_{t,j}=y,\,D_{t,j}^{\mathrm{val}}=\ell\mid\theta_t\right)
= g_D^{\mathrm{rpt}}(d)\,p^y(1-p)^{1-y}\,g_{D,y}^{\mathrm{val}}(\ell),
\end{split}
\end{equation}
where $g_{D,1}^{\mathrm{val}}=g_{D+}^{\mathrm{val}}$ and $g_{D,0}^{\mathrm{val}}=g_{D-}^{\mathrm{val}}$ are defined below. A similar factorization applies to probability mass at $D^{\mathrm{val}}=\infty$. 

We assume the reporting delay has a probability mass function $g_D^{\text{rpt}}$ and denote its distribution function $G_D^{\text{rpt}}$. The survival function is $\bar{G}_D^{\text{rpt}} = 1 - G_D^{\text{rpt}}$. The validation delay is allowed to depend on the latent sign with
\begin{equation}
    g_{D+}^{\text{val}}(\ell) = \Prb(D^{\text{val}} = \ell \mid Y = 1) 
    \quad \text{and} \quad 
    g_{D-}^{\text{val}}(\ell) = \Prb(D^{\text{val}} = \ell \mid Y = 0)
\end{equation}
denoting the conditional probability mass functions for the confirmation and retraction delays for finite $\ell \in \{0,1,2,\dots\}$. Their corresponding distribution functions on the finite delays are $G_{D+}^{\text{val}}$ and $G_{D-}^{\text{val}}$. We define their survival functions by
\begin{equation}
    \bar G_{D+}^{\text{val}}(\ell)
    = \Prb(D^{\text{val}} > \ell \mid Y=1),
    \qquad
    \bar G_{D-}^{\text{val}}(\ell)
    = \Prb(D^{\text{val}} > \ell \mid Y=0),
\end{equation}
so that any probability mass at $D^{\text{val}}=\infty$ is included in the corresponding survival function.

Finally, we define the true-case count as the total number of positive (true) cases at time $t$:
\begin{equation}
    N_t^{+} = \sum\limits_{j = 1}^{M_t} Y_{t,j}
\end{equation}
We denote the expected values of the gross reports and the positive cases as:
\begin{equation}
    \lambda_t^{+} = \E[N_t^{+}] \quad \text{and} \quad  \mu_t = \E[M_t].
\end{equation}
The total number of non-true cases is:
\begin{equation}
    N_t^{-} =  \sum\limits_{j = 1}^{M_t} (1 - Y_{t,j}).
\end{equation}
Because the marks are independent of $M_t$ and $Y_{t,j}\sim\operatorname{Bernoulli}(p)$,
\begin{equation}
    \lambda_t^{+}=p\mu_t,
    \qquad
    \lambda_t^{-}=(1-p)\mu_t.
\end{equation}
We emphasize that $N_t^{-}$ counts non-true cases, not necessarily observed retractions, because a non-true case is allowed to have a validation delay $D^{\mathrm{val}}=\infty$ and therefore may never be retracted.

We assume the distribution of $M_t$ is such that its expected value exists and is finite. 


\subsubsection{The observation map}

We define $\mathcal{O}_{\tau}(X_t)$ as the record generated by a latent point $X_t$ at the now $\tau$, with $\mathcal{O}_{\tau}(X_t)=\emptyset$ denoting a point completely unobserved at the analysis date. Otherwise $\mathcal{O}_{\tau}(X_t)$ contains the observable part of its mark. Define
\begin{equation}
    q_{\emptyset,t,\tau}=\Prb\{\mathcal{O}_{\tau}(X_t)=\emptyset\}.
\end{equation}
For an observable record $x$, let $q_{t,\tau}(x)$ denote its probability mass or density with respect to a dominating measure $\nu_t$ on the observable-record space. Thus,
\begin{equation}
    q_{\emptyset,t,\tau}
    + \int q_{t,\tau}(x)\,\nu_t(dx)
    =1.
\end{equation}
\subsection{The general likelihood}

\begin{theorem}[General censored marked point-process likelihood] 
Suppose that at the now, $\tau$, the observed line-list for event time $t$ contains $k_t$ visible records $\mathcal D_t=\{x_{t,1},\ldots,x_{t,k_t}\}$. Conditional on $\theta_t$, and up to factors independent of $\theta_t$, the likelihood contribution is
\begin{equation}\label{glik}
    \mathcal{L}_t(\theta_t;\mathcal D_t) 
    \propto 
    S_{k_t,t}\big(q_{\emptyset, t, \tau};\theta_t\big) 
    \prod\limits_{i = 1}^{k_t} q_{t,\tau}(x_{t,i}),
\end{equation}
where
\begin{equation}
    S_{k,t}\big(q;\theta_t\big) = \sum\limits_{m=k}^{\infty} \binom{m}{k} q^{m-k} \Prb(M_t = m\mid\theta_t).
\end{equation}
The corresponding unnormalized posterior kernel is $\mathcal K_t(\theta_t;\mathcal D_t)=\pi(\theta_t)\mathcal L_t(\theta_t;\mathcal D_t)$.
\begin{proof}
Conditioning on $M_t=m\geq k_t$ one can choose which $k_t$ of the $m$ latent points produced the visible records. There are $\binom{m}{k_t}$ such choices up to the ordering of the observed rows, whose contribution is independent of $\theta_t$. Each visible record $x_{t,i}$ contributes $q_{t,\tau}(x_{t,i})$. The remaining $m-k_t$ unobserved records must each map to $\emptyset$ with probability $q_{\emptyset, t, \tau}$. Conditional independence of the marks gives
\begin{equation}
    \Prb(x_{t,1},\dots,x_{t,k_t}\mid M_t=m,\theta_t)
    \propto    
    q_{\emptyset,t,\tau}^{m-k_t} \binom{m}{k_t}
    \prod_{i=1}^{k_t}q_{t,\tau}(x_{t,i}).
\end{equation}
Marginalizing over $m$ conditional on $\theta_t$ yields \eqref{glik}.
\end{proof}
\end{theorem}

The following two expressions for $S_{k,t}$ have been derived elsewhere {\color{red} PREVIOUS PAPER}. For $M_t\mid\theta_t\sim\Pois(\mu_t)$, it yields
\begin{equation}\label{poissonS}
    S_{k,t}(q;\theta_t) = \dfrac{\mu_t^k}{k!} \exp\big(-\mu_t(1 - q)\big),
\end{equation}
while for $M_t\mid\theta_t$ following a Negative Binomial distribution parametrized by mean $\mu_t$ and size $\kappa_t$,
\begin{equation}\label{negbinS}
    S_{k,t}(q;\theta_t) 
    = 
    \dfrac{\Gamma(k + \kappa_t)}{\Gamma(\kappa_t)\, k!} 
    \dfrac{\kappa_t^{\kappa_t}\mu_t^k}
    {\big(\kappa_t + \mu_t (1 - q)\big)^{k + \kappa_t}}.
\end{equation}

We note that, given the previous result, in general, for any line-list scenario, we just need to identify the $q_{\cdot,\tau}$ elements. 

\subsection{The general result for line-list data}

For an observed row $j$ for event time $t$ with now $\tau$ that was reported at time $r_j$ denote the observed delay and the observed report's age as:
\begin{equation}
    d_j^{\text{rpt}} = r_j - t, \quad \text{and} \quad a_j = \tau - r_j.
\end{equation}
If the row has been resolved at time $v_j$ we also have the observed validation delay:
\begin{equation}
d_j^{\text{val}} = v_j - r_j.
\end{equation}
In every line-list scenario a latent point has not been observed if its report date has not appeared by $\tau$, therefore the probability of not being observed is the probability of being reported after the now:
\begin{equation}
    q_{\emptyset, t, \tau} 
    = \Prb(D^{\text{rpt}} > d_t^{\star}) 
    = \bar{G}_D^{\text{rpt}}(d_t^{\star}).
\end{equation}
For a row $j$ that has been observed there are three possible scenarios:
\begin{enumerate}
    \item \textbf{Confirmation:} After a delay $d_j^{\text{val}}$ the case has been confirmed as a true case. This happens with probability $p \cdot g_{D+}^{\text{val}}(d_j^{\text{val}})$.
    \item \textbf{Retraction:} After a delay $d_j^{\text{val}}$ the case has been retracted. This happens with probability $(1 -p) \cdot g_{D-}^{\text{val}}(d_j^{\text{val}})$.
    \item \textbf{Pending:} A case at the now $\tau$ is pending, meaning that it has neither been confirmed nor retracted. This happens with probability $p \cdot \bar{G}_{D+}^{\text{val}}(a_j) + (1 - p) \cdot \bar{G}_{D-}^{\text{val}}(a_j)$.
\end{enumerate}
For each row $j$ at time $t$ with now $\tau$, we define the validation-state factor as:
\begin{equation}\label{chilinelist}
    \chi_{t,\tau}(j) = \begin{cases}
        p \cdot g_{D+}^{\text{val}}(d_j^{\text{val}}) & \text{ if row } j \text{ is confirmed,} \\
        (1 -p) \cdot g_{D-}^{\text{val}}(d_j^{\text{val}}) & \text{ if row } j \text{ is retracted,} \\
        p \cdot \bar{G}_{D+}^{\text{val}}(a_j) + (1 - p) \cdot \bar{G}_{D-}^{\text{val}}(a_j) & \text{ if row } j \text{ is pending.} 
    \end{cases}
\end{equation}
We can then rewrite the probability of the observation associated to a specific mark $x_j$ as:
\begin{equation}\label{qlinelist}
    q_{t,\tau}(x_{t,j}) = g_{D}^{\text{rpt}} (d^{\text{rpt}}_j)  \cdot \chi_{t,\tau}(j)
\end{equation}
which combines the probability of the observed reporting delay, $g_{D}^{\text{rpt}} (d^{\text{rpt}}_j)$, with the probability associated with its validation state, $\chi_{t,\tau}(j)$. 

Note that in the case of a model where no confirmations are registered (only retractions) $g_{D+}^{\text{val}}(d) = 0$ for all finite $d$ while the survival function is always equal to one: $\bar{G}_{D+}^{\text{val}}(d) = 1$. Analogously, a model with no retractions (confirmations only) has $g_{D-}^{\text{val}}(d) = 0$ and $\bar{G}_{D-}^{\text{val}}(d) = 1$ for any finite $d$. Finally, a model with no confirmations nor retractions has $\chi_{t,\tau}(j) = 1$ for all $j$, which recovers the ordinary event-to-report observation model of {\color{red} REF}. 

Furthermore, in expressions \eqref{qlinelist} and \eqref{chilinelist} one can substitute $g_{D}^{\text{rpt}}(d)$ or any $g_{D\cdot}^{\text{val}}(d)$ for the probability of $D$ lying on an interval, $\mathcal{I}$, $\Prb(D \in \mathcal{I})$ if the delays are known only approximately. 

We write the likelihood \eqref{glik} associated to the line-list data observed at the now $\tau$ corresponding to an event at time $t$ as:
\begin{equation}\label{linelist}  
    \mathcal{L}_t(\theta_t;\mathcal D_t) 
    \propto 
    S_{k_t,t}\big(\bar{G}_{D}^{\text{rpt}}(d_t^{\star});\theta_t\big) 
    \prod\limits_{j = 1}^{k_t} 
    \big[ g_{D}^{\text{rpt}} (d^{\text{rpt}}_j)  \cdot \chi_{t,\tau}(j)\big],
\end{equation}
where for the Poisson and Negative Binomial case $S_{k_t,t}$ can be written respectively as in \eqref{poissonS} or \eqref{negbinS}. The corresponding unnormalized posterior kernel is $\mathcal K_t(\theta_t;\mathcal D_t)=\pi(\theta_t)\mathcal L_t(\theta_t;\mathcal D_t)$.


\subsection{Count-cumulative data}

In some surveillance systems only cumulative counts are available. In the
cumulative-only observation regime considered here, a report enters the count
when it is recorded and leaves the count only if it is subsequently withdrawn.
Confirmations and latent truth labels are not separately observed. The
observable stream therefore does not identify the probability of truth $p$ and
the conditional retraction-delay distribution $g_{D-}^{\mathrm{val}}$
separately. It identifies their induced, unconditional retraction mechanism.

Accordingly, for this observation regime we use the reduced mark
\begin{equation}\label{reducedcummark}
    X_{t,j}^{\mathrm{cum}}
    =
    \left(D_{t,j}^{\mathrm{rpt}},R_{t,j}\right),
\end{equation}
where $R_{t,j}$ is the age of the report when it is withdrawn:
\begin{equation}
    R_{t,j}
    =
    \begin{cases}
        \ell, & \text{if the report is withdrawn exactly $\ell$ days after it appears},\\
        \infty, & \text{if the report is never withdrawn}.
    \end{cases}
\end{equation}
Conditional on $\theta_t$, the reduced marks are i.i.d., independent of $M_t$,
and satisfy $D^{\mathrm{rpt}}\perp R$. The reporting delay retains probability
mass function $g_D^{\mathrm{rpt}}$. We define the finite-age retraction kernel
by
\begin{equation}\label{hRdef}
    h_R(\ell)
    :=
    \Prb(R=\ell\mid\theta_t),
    \qquad \ell=1,2,\dots.
\end{equation}
This is a possibly defective probability mass function on the finite ages:
\begin{equation}\label{hRmass}
    \sum_{\ell=1}^{\infty}h_R(\ell)\leq 1,
    \qquad
    \pi_\infty
    :=
    \Prb(R=\infty\mid\theta_t)
    =1-\sum_{\ell=1}^{\infty}h_R(\ell).
\end{equation}
The corresponding survival function is
\begin{equation}\label{SRdef}
\begin{split}
    S_R(a)
    &:={}
    \Prb(R>a\mid\theta_t)\\
    &=1-\sum_{\ell=1}^{a}h_R(\ell)
    =\pi_\infty+\sum_{\ell=a+1}^{\infty}h_R(\ell),
    \qquad a=0,1,2,\dots.
\end{split}
\end{equation}
Under the original truth--validation parameterization, these quantities would
be induced by
\begin{equation}\label{collapsedretraction}
    h_R(\ell)
    =(1-p)g_{D-}^{\mathrm{val}}(\ell),
    \qquad
    S_R(a)
    =p+(1-p)\bar G_{D-}^{\mathrm{val}}(a).
\end{equation}
Only the combined quantities on the left-hand side enter the cumulative-data
likelihood. We therefore take $h_R$ (equivalently, $S_R$) as primitive and do
not estimate $p$ and $g_{D-}^{\mathrm{val}}$ separately in this observation
regime.

We impose the following assumptions in the cumulative-only model:
\begin{enumerate}
    \item \textbf{Finite reporting delays:} Every latent event is eventually
    reported, so $D^{\mathrm{rpt}}\in\{0,1,2,\dots\}$ almost surely.
    \item \textbf{Finite age for every observed withdrawal:} A report that is
    eventually withdrawn has $R<\infty$, but no finite upper bound is imposed
    on its retraction age. A report that is never withdrawn has $R=\infty$.
    \item \textbf{No same-day withdrawal:} A report cannot enter and leave the
    count at the same delay. Equivalently, $\Prb(R=0)=0$, which is encoded by
    the support in \eqref{hRdef}. This ensures that every report produces an
    observable positive entry before any later negative update.
\end{enumerate}

For event time $t$, the running total number of unretracted reports by delay
$d$ is
\begin{equation}\label{cumdef}
    C_t(d)
    :=
    \#\left\{
    j:
    D_{t,j}^{\mathrm{rpt}}\leq d,\;
    R_{t,j}>d-D_{t,j}^{\mathrm{rpt}}
    \right\}.
\end{equation}
Thus, a report contributes to $C_t(d)$ after it appears and until it is
withdrawn. The eventually retained count and its expected value are
\begin{equation}\label{retainedcount}
    N_t^{\mathrm{ret}}
    :=
    \sum_{j=1}^{M_t}\mathbf 1\{R_{t,j}=\infty\},
    \qquad
    \lambda_t^{\mathrm{ret}}
    :=
    \E[N_t^{\mathrm{ret}}\mid\theta_t]
    =\pi_\infty\mu_t.
\end{equation}
This is an eventual database-retention estimand, not necessarily a biological
truth estimand. The equality $N_t^{\mathrm{ret}}=N_t^+$ requires the additional
substantive assumption that a report is true if and only if it is never
withdrawn.

The reduction removes the non-identifiable decomposition
$(1-p)g_{D-}^{\mathrm{val}}(\ell)$ from the cumulative likelihood. At a finite
maximum report age $A$, however, the data identify only
$h_R(1),\dots,h_R(A)$ and $S_R(A)$. They cannot separate the reports that will
be withdrawn after age $A$ from the mass $\pi_\infty$ that will never be
withdrawn. Consequently, $N_t^{\mathrm{ret}}$ requires sufficiently long
follow-up or an additional tail restriction, whereas the finite-horizon
observation model below is parameterized entirely by identifiable quantities.

The probability that one latent event has entered the cumulative count by
delay $d$ and remains there at that delay is
\begin{equation}\label{qCdef}
    q_C(d)
    :=
    \Prb\left(
    D^{\mathrm{rpt}}\leq d,\;
    R>d-D^{\mathrm{rpt}}
    \mid\theta_t
    \right)
    =
    \sum_{r=0}^{d}
    g_D^{\mathrm{rpt}}(r)S_R(d-r).
\end{equation}
It follows for any epidemic process with finite mean that
\begin{equation}\label{cummean}
    \E[C_t(d)\mid\theta_t]
    =
    \mu_t q_C(d).
\end{equation}

\begin{proposition}[Poisson marginal for a cumulative level]
\label{prop:poissoncumulative}
If $M_t\mid\theta_t\sim\Pois(\mu_t)$, then for every $d\geq0$,
\begin{equation}\label{poissoncumulative}
    C_t(d)\mid\theta_t
    \sim
    \Pois\left(
    \mu_t
    \sum_{r=0}^{d}
    g_D^{\mathrm{rpt}}(r)S_R(d-r)
    \right).
\end{equation}
\begin{proof}
Each of the $M_t$ latent events belongs to the count at delay $d$ independently
with probability $q_C(d)$. The result follows by Poisson thinning.
\end{proof}
\end{proposition}

Proposition~\ref{prop:poissoncumulative} is a one-delay marginal statement.
The levels $C_t(0),C_t(1),\dots$ are dependent because the same report may
contribute at several delays. Hence, a product of the Poisson masses in
\eqref{poissoncumulative} over $d$ is a composite likelihood rather than the
exact joint likelihood. Notice also that a sequence of cumulative levels need
not be monotone: newly appearing reports increase it, while withdrawals can
decrease it.

For our model, we define the signed updates
\begin{equation}\label{signedupdates}
    \Delta^0_t = C_t(0),
    \qquad
    \Delta^d_t = C_t(d) - C_t(d-1), \qquad d\geq 1,
\end{equation}
where a positive $\Delta^d_t$ represents net additions and a negative value represents net withdrawals.

We will rewrite the updates $\{\Delta^d_t\}$ in terms of latent report trajectories. Fix the now ($\tau$) and write $d_t^{\star}=\tau-t$. At this horizon, every latent event that happened at time $t$ belongs to exactly one of the following classes:
\begin{itemize}
    \item \textbf{Not yet reported.} Define
    \begin{equation}
        Z_{t,\tau}
        :=
        \#\{j:D^{\mathrm{rpt}}_{t,j}>d_t^{\star}\}.
    \end{equation}
    Its trajectory probability is
    \begin{equation}\label{qemptycum}
        \tilde q_{\emptyset,t,\tau}
        =
        \bar G_D^{\mathrm{rpt}}(d_t^{\star}).
    \end{equation}

    \item \textbf{Reported at delay $d$ and still unretracted by the now.} For $0\leq d\leq d_t^{\star}$ define
    \begin{equation}\label{Udef}
        U_{t,\tau}^{d}
        :=
        \#\left\{
        j:
        D^{\mathrm{rpt}}_{t,j}=d,
        \text{ and case }j\text{ is unretracted at }\tau
        \right\}.
    \end{equation}
    Its trajectory probability is
    \begin{equation}\label{qUcum}
        \tilde q_{t,\tau}^{U}(d)
        =
        g_D^{\mathrm{rpt}}(d)
        S_R(d_t^{\star}-d).
    \end{equation}

    \item \textbf{Reported at delay $d_1$ and retracted at a later delay $d_2$.} For $0\leq d_1<d_2\leq d_t^{\star}$ define
    \begin{equation}\label{Wdef}
        W_{t,\tau}^{d_1,d_2}
        :=
        \#\left\{
        j:
        D^{\mathrm{rpt}}_{t,j}=d_1,
        D^{\mathrm{rpt}}_{t,j}+R_{t,j}=d_2
        \right\}.
    \end{equation}
    Its trajectory probability is
    \begin{equation}\label{qWcum}
        \tilde q_{t,\tau}^{W}(d_1,d_2)
        =
        g_D^{\mathrm{rpt}}(d_1)
        h_R(d_2-d_1).
    \end{equation}
\end{itemize}
The probabilities of the trajectory classes satisfy
\begin{equation}\label{trajectorypartition}
    \tilde q_{\emptyset,t,\tau}
    +
    \sum_{d=0}^{d_t^{\star}}\tilde q_{t,\tau}^{U}(d)
    +
    \sum_{0\leq d_1<d_2\leq d_t^{\star}}
    \tilde q_{t,\tau}^{W}(d_1,d_2)
    =1.
\end{equation}
This follows because the three classes are mutually exclusive and exhaustive at the analysis time $\tau$.

Using these trajectory counts, for every $d=0,\dots,d_t^{\star}$,
\begin{equation}\label{deltadecomposition}
    \Delta_t^d
    =
    U_{t,\tau}^{d}
    +
    \sum_{d_2=d+1}^{d_t^{\star}}W_{t,\tau}^{d,d_2}
    -
    \sum_{d_1=0}^{d-1}W_{t,\tau}^{d_1,d},
\end{equation}
where an empty sum is defined to be zero. The first term counts reports appearing at delay $d$ that remain unretracted at the now; the second counts reports appearing at delay $d$ that were subsequently retracted by the now and therefore contributed $+1$ when first reported; the final term counts earlier reports that were retracted at delay $d$ and therefore contribute $-1$ to that update.

\subsubsection{Marginal distributions at a fixed $d$}

For a fixed delay $d$, define the total number of additions and withdrawals at that delay as
\begin{equation}\label{Adef}
    \mathcal A_t^d
    :=
    U_{t,\tau}^{d}
    +
    \sum_{d_2=d+1}^{d_t^{\star}}W_{t,\tau}^{d,d_2}, \quad
\text{and}\quad
    \mathcal W_t^d
    :=
    \sum_{d_1=0}^{d-1}W_{t,\tau}^{d_1,d}.
\end{equation}
Then we can write the updates as:
\begin{equation}\label{deltaAB}
    \Delta_t^d=\mathcal A_t^d-\mathcal W_t^d.
\end{equation}
We remark that the trajectory classes contributing to $\mathcal A_t^d$ and $\mathcal W_t^d$ are disjoint.

\paragraph{Poisson marginal}
Suppose, conditional on $\theta_t$, the epidemic process follows a Poisson distribution:
\begin{equation}
    M_t\mid\theta_t\sim\Pois(\mu_t).
\end{equation}
By Poisson thinning, the counts associated with mutually exclusive trajectory classes are independent Poisson random variables. Therefore $\mathcal A_t^d$ and $\mathcal W_t^d$ are independent Poisson random variables. Define their means as
\begin{equation}\label{alphadef}
    \alpha_t^d
    :=
    \E[\mathcal A_t^d\mid\theta_t]=
    \mu_t\left[
    \tilde q_{t,\tau}^{U}(d)
    +
    \sum_{d_2=d+1}^{d_t^{\star}}
    \tilde q_{t,\tau}^{W}(d,d_2)
    \right] = \mu_t g_D^{\mathrm{rpt}}(d)
\end{equation}
and 
\begin{equation}\label{omegadef}
    \omega_t^d
    :=
    \E[\mathcal W_t^d\mid\theta_t]
    =
    \mu_t
    \sum_{d_1=0}^{d-1}
    \tilde q_{t,\tau}^{W}(d_1,d)
    =
    \mu_t
    \sum_{d_1=0}^{d-1}
    g_D^{\mathrm{rpt}}(d_1)
    h_R(d-d_1).
\end{equation}

The equality for $\alpha_t^d$ follows because $\mathcal A_t^d$ counts every
event whose report first appears at delay $d$, irrespective of whether it is
retracted later. Its mean is therefore the expected number of events multiplied
by the probability of first appearing at delay $d$.


\begin{proposition}[Skellam marginal under a Poisson epidemic process]\label{prop:poissonskellam}
Conditional on $\theta_t$, for every $d=0,\dots,d_t^{\star}$,
\begin{equation}\label{poissonskellam}
    \Delta_t^d\mid\theta_t
    \sim
    \Sk(\alpha_t^d,\omega_t^d),
\end{equation}
where $\alpha_t^d$ and $\omega_t^d$ are defined in \eqref{alphadef} and \eqref{omegadef}. 
\begin{proof}
Under Poisson thinning, all $U$ and $W$ trajectory counts are mutually independent Poisson variables with means equal to $\mu_t$ times their trajectory probabilities. The sum defining $\mathcal A_t^d$ is therefore Poisson with mean $\alpha_t^d$, and the sum defining $\mathcal W_t^d$ is Poisson with mean $\omega_t^d$. The two sums involve disjoint trajectory classes and are therefore independent. By \eqref{deltaAB}, $\Delta_t^d$ is the difference of two independent Poisson variables, which is a Skellam random variable with parameters $(\alpha_t^d,\omega_t^d)$. 
\end{proof}
\end{proposition}

We remind the readers that for $\alpha>0$ and $\omega>0$, the Skellam probability mass function is
\begin{equation}\label{skellampmf}
    \Prb\{\Sk(\alpha,\omega)=z\}
    =
    \exp[-(\alpha+\omega)]
    \left(\frac{\alpha}{\omega}\right)^{z/2}
    I_{|z|}\left(2\sqrt{\alpha\omega}\right),
    \qquad z\in\mathbb Z,
\end{equation}
where $I_\nu$ denotes the modified Bessel function of the first kind. Boundary cases are defined by continuity. If $\omega=0$, the distribution is $\Pois(\alpha)$ on the non-negative integers; if $\alpha=0$, then $-\Delta_t^d$ is $\Pois(\omega)$ and $\Delta_t^d$ is supported on the non-positive integers.



\paragraph{Negative binomial marginal} If $M_t$ follows a Negative Binomial distribution with mean $\mu_t$ and size $\kappa_t$ we can utilize the standard Gamma--Poisson representation of the Negative Binomial to obtain the marginal distribution at a delay $d$. 

Let
\begin{equation}\label{xidef}
\Xi_t\sim\operatorname{Gamma}(\kappa_t,\kappa_t),
\end{equation}
where we follow a shape--rate parametrization such that
\begin{equation}
    \E[\Xi_t]=1,
    \qquad
    \Var(\Xi_t)=\frac{1}{\kappa_t}.
\end{equation}
Then
\begin{equation}\label{nbmixture}
    M_t\mid\Xi_t=\xi
    \sim
    \Pois(\xi\mu_t)
\end{equation}
has the marginal distribution $M_t\sim\NB(\mu_t,\kappa_t)$.

Conditional on $\Xi_t=\xi$, the Poisson splitting argument also applies here with every trajectory being multiplied by $\xi$. In particular,
\begin{equation}
    \mathcal{A}_t^d\mid\Xi_t=\xi
    \sim
    \Pois(\xi\alpha_t^d),
    \qquad
    \mathcal{W}_t^d\mid\Xi_t=\xi
    \sim
    \Pois(\xi\omega_t^d),
\end{equation}
with the two variables being conditionally independent.

\begin{proposition}[Gamma-mixed Skellam marginal under Negative Binomial epidemic process]\label{prop:nbskellam}
For every $d=0,\dots,d_t^{\star}$,
\begin{equation}\label{conditionalnbskellam}
    \Delta_t^d\mid\Xi_t=\xi
    \sim
    \Sk(\xi\alpha_t^d,\xi\omega_t^d).
\end{equation}
Consequently, the marginal probability mass function of $\Delta_t^d$ is the Gamma mixture
\begin{equation}\label{gammamixedskellam}
    \Prb(\Delta_t^d=m)
    =
    \int_0^\infty
    \Prb\{\Sk(\xi\alpha_t^d,\xi\omega_t^d)=m\}
    f_{\Xi_t}(\xi)\,d\xi,
    \qquad m\in\mathbb Z.
\end{equation}
where $f_{\Xi_t}(\xi)$ represents the Gamma density function with shape and rate $\kappa_t$.
\end{proposition}


As a sidenote, the results in Propositions~\ref{prop:poissonskellam} and \ref{prop:nbskellam} are one-delay marginal distributions. Updates at different delays are generally dependent because a single trajectory $W_{t,\tau}^{d_1,d_2}$ contributes $+1$ at $d_1$ and $-1$ at $d_2$. Consequently, multiplying the delay-specific Skellam or Gamma-mixed Skellam masses across $d$ gives a composite likelihood, not the exact joint likelihood for the full update vector.

\subsubsection{An overdispersed, zero-inflated marginal}
\label{sec:ztnb}

Propositions~\ref{prop:poissonskellam} and~\ref{prop:nbskellam} tie the
dispersion of an update to its mean. A $\Sk(\alpha,\omega)$ variable has
\begin{equation}
    \E[\Delta_t^d\mid\theta_t]=\alpha_t^d-\omega_t^d,
    \qquad
    \Var(\Delta_t^d\mid\theta_t)=\alpha_t^d+\omega_t^d,
\end{equation}
so the only way to place appreciable probability on large updates is to make
$\alpha_t^d+\omega_t^d$ large, which simultaneously drives
$\Prb(\Delta_t^d=0)$ toward zero. In surveillance streams where the cumulative
count is essentially complete at first publication, the opposite pattern is
typical: the great majority of updates are exactly null, and the remainder are
occasionally very large. A gamma mixture as in \eqref{gammamixedskellam} does not
resolve this, because it scales $\alpha_t^d$ and $\omega_t^d$ together and
therefore inflates the tail and depletes the atom at the same time.

We therefore introduce a marginal that separates three features of an update
which the Skellam family forces to share a single parameter: whether the update
is null at all, its direction, and its magnitude.

\begin{definition}[Zero-truncated negative binomial]
For a parent mean $m>0$ and size $\varsigma>0$, write
\begin{equation}\label{ztnbpmf}
    \Prb\{\operatorname{ZTNB}^{\mathrm{par}}(m,\varsigma)=k\}
    =
    \frac{\Prb\{\NB(m,\varsigma)=k\}}
         {1-\Prb\{\NB(m,\varsigma)=0\}},
    \qquad k=1,2,\dots,
\end{equation}
where $\NB(m,\varsigma)$ has mean $m$ and size $\varsigma$. Its mean is
\emph{not} $m$:
\begin{equation}\label{ztnbmean}
    \E\left[\operatorname{ZTNB}^{\mathrm{par}}(m,\varsigma)\right]
    =
    \frac{m}{1-\Prb\{\NB(m,\varsigma)=0\}}
    =
    \frac{m}{1-\left(\frac{\varsigma}{\varsigma+m}\right)^{\varsigma}}
    =:
    \Psi_\varsigma(m).
\end{equation}
\end{definition}

\begin{lemma}[Mean reparameterization]\label{lem:ztnbmean}
For every $\varsigma>0$ the map $\Psi_\varsigma$ of \eqref{ztnbmean} is a
continuous, strictly increasing bijection from $(0,\infty)$ onto $(1,\infty)$.
\begin{proof}
Continuity and strict monotonicity are immediate from
$\Psi_\varsigma(m)=m/(1-P_0(m))$ with $P_0(m)=(\varsigma/(\varsigma+m))^\varsigma$
strictly decreasing in $m$, so that both $m$ and $1/(1-P_0(m))$ are strictly
increasing and positive. As $m\downarrow0$, the zero-truncated law concentrates
on $\{1\}$ and $\Psi_\varsigma(m)\downarrow1$; as $m\uparrow\infty$,
$P_0(m)\downarrow0$ and $\Psi_\varsigma(m)\sim m\uparrow\infty$.
\end{proof}
\end{lemma}

Lemma~\ref{lem:ztnbmean} lets us index the family by its \emph{own} mean. For
$\nu>1$ write $\operatorname{ZTNB}(\nu,\varsigma)
:=\operatorname{ZTNB}^{\mathrm{par}}(\Psi_\varsigma^{-1}(\nu),\varsigma)$, so that
$\E[\operatorname{ZTNB}(\nu,\varsigma)]=\nu$ exactly. The constraint $\nu>1$ is
not a restriction but a necessity: the law is supported on $\{1,2,\dots\}$.

\begin{definition}[Hurdle update model]\label{def:hurdle}
Let $\varsigma_R>0$ and let $\pi_t^d\in(0,1]$ denote the probability that update
$d$ at event time $t$ is non-null, subject to the admissibility condition
\begin{equation}\label{admissible}
    \pi_t^d\leq\alpha_t^d+\omega_t^d .
\end{equation}
Given $\alpha_t^d,\omega_t^d$ from \eqref{alphadef} and \eqref{omegadef} with
$\alpha_t^d+\omega_t^d>0$, define
\begin{equation}\label{hurdletheta}
    \vartheta_t^d
    :=
    \frac{\alpha_t^d}{\alpha_t^d+\omega_t^d},
    \qquad
    \nu_t^d
    :=
    \frac{\alpha_t^d+\omega_t^d}{\pi_t^d}\;(\geq1\text{ by }\eqref{admissible}),
\end{equation}
and let
\begin{equation}\label{hurdlelaw}
    \Delta_t^d
    =
    \begin{cases}
        0, & \text{with probability } 1-\pi_t^d,\\
        +M_t^d, & \text{with probability } \pi_t^d\vartheta_t^d,\\
        -M_t^d, & \text{with probability } \pi_t^d(1-\vartheta_t^d),
    \end{cases}
    \qquad
    M_t^d\sim\operatorname{ZTNB}(\nu_t^d,\varsigma_R).
\end{equation}
\end{definition}

Condition \eqref{admissible} has a direct mechanistic reading: since every
non-null update has magnitude at least one, the expected absolute flow
$\alpha_t^d+\omega_t^d$ bounds how often an update can be non-null at all. One
cannot have more cohorts moving than there are expected report and retraction
events.

The direction $\vartheta_t^d$ and the magnitude scale $\nu_t^d$ are not free:
both are determined by the same $\alpha_t^d$ and $\omega_t^d$ that the reporting
and retraction mechanism already supplies. Only $\pi_t^d$ and $\varsigma_R$ are
additional.

\begin{proposition}[Mean preservation]\label{prop:hurdlemean}
Under Definition~\ref{def:hurdle}, for every admissible $\pi_t^d$,
\begin{equation}\label{hurdlemean}
    \E[\Delta_t^d\mid\theta_t]
    =
    \alpha_t^d-\omega_t^d,
\end{equation}
and consequently
\begin{equation}\label{hurdlecummean}
    \E[C_t(d)\mid\theta_t]
    =
    \mu_t q_C(d)
    =
    \mu_t\sum_{r=0}^{d}g_D^{\mathrm{rpt}}(r)S_R(d-r).
\end{equation}
\begin{proof}
Write $\alpha=\alpha_t^d$, $\omega=\omega_t^d$, $\pi=\pi_t^d$. By
Lemma~\ref{lem:ztnbmean} and the indexing convention following it,
$\E[M_t^d]=\nu_t^d=(\alpha+\omega)/\pi$ exactly. Conditioning on whether the
update is null,
\begin{equation}
\begin{split}
    \E[\Delta_t^d\mid\theta_t]
    &=
    \pi\vartheta_t^d\,\E[M_t^d]
    -
    \pi\left(1-\vartheta_t^d\right)\E[M_t^d]
    =
    \pi\,\E[M_t^d]\left(2\vartheta_t^d-1\right)\\
    &=
    \pi\cdot\frac{\alpha+\omega}{\pi}
    \cdot
    \left(\frac{2\alpha}{\alpha+\omega}-1\right)
    =
    (\alpha+\omega)\cdot\frac{\alpha-\omega}{\alpha+\omega}
    =
    \alpha-\omega,
\end{split}
\end{equation}
and the factor $\pi$ cancels. Summing over $d$ and telescoping
$\sum_{e\leq d}\Delta_t^e=C_t(d)$ gives \eqref{hurdlecummean}, since
$\sum_{e\leq d}(\alpha_t^e-\omega_t^e)=\mu_t q_C(d)$ by \eqref{cummean}.
\end{proof}
\end{proposition}

\begin{remark}
The indexing of Lemma~\ref{lem:ztnbmean} is essential and easy to get wrong. If
$(\alpha+\omega)/\pi$ is passed as the \emph{parent} mean $m$ rather than as the
mean $\nu$, then $\E[M_t^d]=\Psi_{\varsigma_R}(m)>m$ and the realized mean
becomes $(\alpha-\omega)/\left(1-P_0\right)$ rather than $\alpha-\omega$. The
discrepancy is severe at small $\varsigma_R$: at $\varsigma_R=0.015$ the
inflation factor $1/(1-P_0)$ ranges from roughly $8$ to $37$ over the parameter
values encountered in practice, so \eqref{hurdlecummean} fails badly and the
identification argument no longer transfers.
\end{remark}

Proposition~\ref{prop:hurdlemean} is the reason for the particular
parameterization in \eqref{hurdletheta}. The first moment of the update, and
hence of the cumulative level, is exactly that of the Skellam model of
Proposition~\ref{prop:poissonskellam}, for any choice of $\pi_t^d$. The hurdle
therefore changes only the dispersion and the atom at zero; every statement about
what the cumulative data do and do not identify --- in particular that they pin
$h_R(1),\dots,h_R(A)$ and $S_R(A)$ but not $\pi_\infty$, and that $p$ and
$g_{D-}^{\mathrm{val}}$ are not separately identified --- carries over unchanged.

The variance is inflated relative to the Skellam. Writing
$\varsigma_R^{-1}$ for the negative binomial dispersion, a direct computation
gives
\begin{equation}\label{hurdlevar}
    \Var(\Delta_t^d\mid\theta_t)
    =
    \pi_t^d\,\E\left[(M_t^d)^2\right]
    -
    \left(\alpha_t^d-\omega_t^d\right)^2,
\end{equation}
which grows as $\pi_t^d$ decreases at fixed $\alpha_t^d+\omega_t^d$: the same
expected flow is delivered by rarer but larger updates. The limit
$\pi_t^d\to1$ with $\varsigma_R\to\infty$ recovers a law with the Skellam's mean
and comparable dispersion, so \eqref{hurdlelaw} nests the behaviour of
Proposition~\ref{prop:poissonskellam} in that corner.

In applications $\pi_t^d$ is naturally modelled as a function of the report age
$d$, since the propensity of a cohort to move at all declines sharply as it
matures, and may also depend on whether the preceding update was non-null, which
introduces a two-state dependence across consecutive delays. Note that a product
over $d$ of the masses implied by \eqref{hurdlelaw} remains a composite
likelihood for the same reason as before: it does not encode the matched
report--retraction trajectories that link $d_1$ and $d_2$.

\subsubsection{The general likelihood for count-cumulative data}

Let
\begin{equation}
    \boldsymbol{\Delta}_t
    =
    \left(
    \Delta_t^0,\Delta_t^1,\dots,\Delta_t^{d_t^{\star}}
    \right)
\end{equation}
denote the observed signed-update vector, and let
\begin{equation}
    \boldsymbol{\delta}_t
    = \left(\delta^0_t,\delta^1_t,\dots,\delta^{d_t^{\star}}_t\right)
\end{equation}
denote one possible observed value. Define $\mathfrak C_{t,\tau}(\boldsymbol\delta_t)$ as the collection of all non-negative integer arrays
\begin{equation}
    \left(
    \{u_d\}_{d=0}^{d_t^{\star}},
    \{w_{d_1,d_2}\}_{0\leq d_1<d_2\leq d_t^{\star}}
    \right)
\end{equation}
that satisfy
\begin{equation}\label{compatibility}
    \delta^d_t
    =
    u_d
    +
    \sum_{d_2=d+1}^{d_t^{\star}}w_{d,d_2}
    -
    \sum_{d_1=0}^{d-1}w_{d_1,d},
    \qquad
    d=0,\dots,d_t^{\star},
\end{equation}
with empty sums defined as zero. Thus $\mathfrak C_{t,\tau}(\boldsymbol\delta_t)$ contains every possible collection of individual report and retraction trajectories that produces the observed vector $\boldsymbol\delta_t$.

For a particular compatible allocation $(\boldsymbol u,\boldsymbol w)$ define
\begin{equation}\label{Kcum}
    K(\boldsymbol u,\boldsymbol w)
    =
    \sum_{d=0}^{d_t^{\star}}u_d
    +
    \sum_{0\leq d_1<d_2\leq d_t^{\star}}w_{d_1,d_2}.
\end{equation}
This is the number of events that have been reported by the now under that trajectory allocation.

\begin{theorem}[General count-cumulative likelihood]\label{thm:generalcum}
Conditional on $\theta_t$, suppose that the observed count-cumulative data at now $\tau$ for event time $t$ are summarized by $\boldsymbol\Delta_t=\boldsymbol\delta_t$. Then the exact likelihood contribution is
\begin{equation}\label{generalcumlik}
\begin{split}
    \mathcal L_t^{\mathrm{cum}}(\theta_t;\boldsymbol\delta_t)
    &:={}
    \Prb(\boldsymbol\Delta_t=\boldsymbol\delta_t\mid\theta_t)\\
    &=
    \sum_{(\boldsymbol u,\boldsymbol w)\in
    \mathfrak C_{t,\tau}(\boldsymbol\delta_t)}
    \frac{K(\boldsymbol u,\boldsymbol w)!}
    {
    \displaystyle\prod_{d=0}^{d_t^{\star}}u_d!
    \displaystyle\prod_{0\leq d_1<d_2\leq d_t^{\star}}w_{d_1,d_2}!}
    \,
    S_{K(\boldsymbol u,\boldsymbol w),t}
    \left(\tilde q_{\emptyset,t,\tau};\theta_t\right)\\
    &\qquad\times
    \prod_{d=0}^{d_t^{\star}}
    \left[\tilde q_{t,\tau}^{U}(d)\right]^{u_d}
    \prod_{0\leq d_1<d_2\leq d_t^{\star}}
    \left[\tilde q_{t,\tau}^{W}(d_1,d_2)\right]^{w_{d_1,d_2}}.
\end{split}
\end{equation}
If $\mathfrak C_{t,\tau}(\boldsymbol\delta_t)$ is empty, the likelihood is zero. 

\begin{proof}
Fix one compatible trajectory allocation $(\boldsymbol u,\boldsymbol w)\in\mathfrak C_{t,\tau}(\boldsymbol\delta_t)$ and write
\begin{equation}
    K=K(\boldsymbol u,\boldsymbol w).
\end{equation}
Conditional on $(M_t=m,\theta_t)$ with $m\geq K$, exactly $K$ latent events have been reported by the now and $m-K$ have not yet been reported. By conditional independence of the marks and the trajectory partition \eqref{trajectorypartition}, the vector consisting of the unreported count, the $U$ counts, and the $W$ counts is multinomial:
\begin{equation}
\begin{split}
    &\Prb\left(
    Z_{t,\tau}=m-K,
    \{U_{t,\tau}^{d}=u_d\},
    \{W_{t,\tau}^{d_1,d_2}=w_{d_1,d_2}\}
    \mid M_t=m,\theta_t
    \right)\\
    &\quad=
    \frac{m!}
    {(m-K)!
    \displaystyle\prod_d u_d!
    \displaystyle\prod_{d_1<d_2}w_{d_1,d_2}!}
    \left(\tilde q_{\emptyset,t,\tau}\right)^{m-K}
    \prod_d
    \left[\tilde q_{t,\tau}^{U}(d)\right]^{u_d}
    \prod_{d_1<d_2}
    \left[\tilde q_{t,\tau}^{W}(d_1,d_2)\right]^{w_{d_1,d_2}}.
\end{split}
\end{equation}
Using
\begin{equation}
    \frac{m!}{(m-K)!}
    =
    K!\binom{m}{K},
\end{equation}
and marginalizing over $M_t$ conditional on $\theta_t$ gives the probability of that particular trajectory allocation:
\begin{equation}
\begin{split}
    &\frac{K!}
    {\displaystyle\prod_d u_d!
    \displaystyle\prod_{d_1<d_2}w_{d_1,d_2}!}
    S_{K,t}\left(\tilde q_{\emptyset,t,\tau};\theta_t\right)
    \prod_d
    \left[\tilde q_{t,\tau}^{U}(d)\right]^{u_d}
    \prod_{d_1<d_2}
    \left[\tilde q_{t,\tau}^{W}(d_1,d_2)\right]^{w_{d_1,d_2}}.
\end{split}
\end{equation}
The observed vector $\boldsymbol\delta_t$ does not identify which compatible trajectory allocation occurred. Because the allocations in $\mathfrak C_{t,\tau}(\boldsymbol\delta_t)$ are mutually exclusive, their probabilities must be summed. This yields \eqref{generalcumlik}. 
\end{proof}
\end{theorem}

Direct evaluation of \eqref{generalcumlik} requires enumeration of all compatible trajectory allocations in $\mathfrak C_{t,\tau}(\boldsymbol\delta_t)$ and may become computationally prohibitive as the observable horizon and counts increase. The one-delay marginal models above avoid this enumeration, but a product of their marginal masses should be interpreted as a composite likelihood because it does not retain the cross-delay dependence induced by matched report--retraction trajectories.

\end{document}
